\documentclass[12pt]{article}

\usepackage[utf8]{inputenc}
\usepackage[spanish]{babel}

\usepackage{amssymb, amsmath}

\usepackage{geometry}
\geometry{margin=2.54cm}

\pagestyle{plain}

\begin{document}

% Portada
\begin{titlepage}
	\centering
	{\scshape\LARGE Universidad Aztlán \par}
	\vspace{1cm}
	{\scshape\Large Campus Playa del Carmen \par}
	\vspace{3cm}
	{\scshape\Huge Álgebra de Vectores \par}
	\vspace{1cm}
	{\itshape\Large Introducción \par}
	\vspace{3cm}
	{\Large José Rosendo Camal Teh \par}
	\vfill
	{\Large Ingeniería en Sistemas Computacionales \par}
	{\Large Tercer Cuatrimestre \par}
	{\Large Cálculo Vectorial \par}
	\vfill
	{\Large Prof. Luis Alberto Garibaldi Leciwr \par}
	\vfill
	{\Large 13 de mayo de 2026 \par}
	\vfill
	{\Large Playa del Carmen, Quintana Roo \par}
\end{titlepage}

\section*{Actividades}

% Actividad 1
\subsection*{Actividad 1}

\textbf{Instrucciones:} Encontrar la magnitud de:
\[ \vec{A} = \left( 5, 12 \right) \]

\textbf{Solución:}
La magnitud del vector $\vec{A}$ es:

\begin{align*}
	|\vec{A}| &= \sqrt{ (5)^{2} + (12)^2 } \\
	&= \sqrt{ 25 + 144 } \\
	&= \sqrt{169} \\
	&= \sqrt{ 13^{2} } \\
	&= 13 \\
	|\vec{A}| &= 13
\end{align*}
\newline

% Actividad 2
\subsection*{Actividad 2}

\textbf{Instrucciones:} Realiza la suma:
\begin{align*}
	\vec{A} = \left( 2, 4 \right) \\
	\vec{B} = \left( 6, 1 \right) 
\end{align*}

\textbf{Solución:}
\begin{align*}
	\vec{A} + \vec{B} &= \left( 2 + 6, 4 + 1 \right) \\
	&= \left( 8, 5 \right) \\
	\vec{A} + \vec{B} &= \left( 8, 5 \right)
\end{align*}
\newline


% Actividad 3
\subsection*{Actividad 3}

\textbf{Instrucciones:} Un montacargas se mueve:
\begin{itemize}
	\item 7 m al este
	\item 4 m al norte
\end{itemize}

Representa el vector y calcula su magnitud.
\newline

\textbf{Solución:}
El vector del movimiento del montacargas está representado por el movimiento horizontal y vertical en el plano cartesiano. En analogía, el movimiento horizontal ---representado en el eje x--- corresponde a la traslación de los puntos cardinales de este a oeste, o vicerversa. Mientras que el movimiento vertical ---representado en el eje y--- corresponde a la traslación de norte a sur (o viceversa). En tal caso, el movimiento horizontal es 7 y el vertical es 4, y por lo tanto, el vector se representa en la forma siguiente:

\[ \vec{M} = \left( 7, 4 \right) \]

La magnitud del vector $\vec{M}$ es igual a la raíz cuadrada de la suma de los cuadrados de sus componentes ---aplicación del Teorema de Pitágoras.

\begin{align*}
	|\vec{M}| &= \sqrt{ (7)^2 + (4)^2 } \\
	&= \sqrt{ 49 + 16 } \\
	&= \sqrt{65} \\
	|\vec{M}| &= \sqrt{65}
\end{align*}

De esta manera, $|\vec{M}| = \sqrt{65} \approx$ 8.062.
\end{document}
